\section*{Limitations}
\label{sec:limitations}

\paragraph{Scope.}
All layer-level and probing results come from OLMo-2-7B. We replicate the
lookahead curves on OLMo-3-7B and generation on Qwen2.5-14B
(Appendix~\ref{app:details}) and the qualitative pattern holds, but we have not
shown that the layer localisation or the probe's profile transfers. The
Qwen2.5-14B result is worth flagging on its own: at roughly twice the parameters
it does not clearly beat OLMo-2-7B here, so recovering phrases this way does not
look like a simple function of general model capability. The phrases are a
deliberately favourable case, and a narrow one: frequent fixed English
expressions, unbalanced across categories (idioms are $82\%$ of pooled trials),
with the landmark and film-title lists taken from public web pages rather than
curated resources. Nothing here speaks to compositional continuations, to rarer
expressions, or to other languages.

\paragraph{What a high recovery rate does not settle.}
Autoregressive models are trained to continue text, so a high recovery rate on
its own does not isolate anything specific to fixed expressions. We did not run
a control of frequent compositional spans matched for next-token predictability,
so we cannot say whether these rates exceed what equally predictable ordinary
continuations would give. Nor did we count how often these phrases occur in
pretraining, although we chose OLMo-2 partly because its training data is
public. Frequency is a plausible common cause of several results at once: the
spread between phrases in \S\ref{sec:stability}, the category differences, and
which phrases become recoverable early. A frequency-matched control and a
correlation against corpus counts are the two analyses that would do most to
separate phrase identity from exposure, and both remain to be done. The success
criterion adds slack in both directions: substring containment credits any
output that contains the target without asking why it appeared, and penalises
accepted variants of the same expression (\S\ref{sec:falsepos}).

\paragraph{How firmly the numbers should be held.}
Our primary Patchscopes measurement unions eight of the model's $32$ layers, so
that every per-layer, per-category and probing result rests on the same stored
states. That undercounts a full sweep by six to eight points and is the single
largest contributor to the gap between the two measurements
(\S\ref{sec:disagreement}). Where the $32$-layer sweep is available we report it
alongside, but it is not available for the per-category or probing analyses,
which should be read as lower bounds on Patchscopes. We report no confidence
intervals. The probe numbers come from one phrase-level split at seed $42$, and
$44.8\%$ of the collected contexts are duplicates (\S\ref{sec:setup}), so the
effective sample behind every instance-level number is closer to $5{,}498$ than
to $9{,}955$ and any interval computed from the instance count would be too
narrow. The right unit for resampling is the phrase rather than the instance,
and a phrase-clustered bootstrap over the central comparisons is the obvious
next step. Differences of one or two points, such as those between models in
Appendix~\ref{app:details}, should not be read as meaningful. Finally, we
implement no decoding scheme and report no latency or throughput measurement:
whether a probe of this kind is cheap enough to pay for itself is open.

\section*{Ethics Statement}

This work analyses publicly available language models on publicly available
text. The phrase sets are drawn from published idiom corpora and public lists of
landmarks and film titles; the contexts are sampled from FineWeb-Edu, a filtered
web corpus whose documents we did not inspect individually and which may contain
the biases of web text. We do not release any model and the analysis does not involve human subjects.
We did not set out to collect or analyse personal information, but the
underlying web corpus was not inspected document by document and may contain
incidental personal data. The efficiency direction this work
motivates, if pursued, would reduce inference cost; we make no claim to have
demonstrated that reduction here.
